% Page budget: 1.55 pages | Owns: augmentation topology, additional states, encoder, and closed-loop objective.
% WRITING GUIDE
% Purpose: Define the final augmentation architecture and explain how it is trained inside the known feedback loop.
% Start here: Current implementation decisions D-129, D-140, D-141, D-180, D-220, D-233, and D-237 in docs/decisions.md.
% Supporting material: Quinten's 18 September outline feedback, the Meeting-audit synthesis, Thesis-writeup/Documentation/TRAINING-DESIGN.md, and docs/writeup figures.
% Mandatory papers: Read Hoekstra's 2025 LFR augmentation paper, the 2026 first-principles augmentation paper, the encoder-initialization paper, and Kessels' Chapter 5 before writing this section.
% Research-plan anchor: Preserve Aspect 2's dynamic parallel augmentation objective; document the departure from open-loop training and earlier state routing.
% Current decisions: Separate physical and additional states, keep the controller outside the model, score the complete training window without burn-in, and use one network for the physical correction and added-state update.
% Choices to justify: Final reported routing, added-state dimension, ANN width and depth, zero initialization, encoder history, closed-loop objective, rollout length, full-window scoring, and validation horizon. Use the configuration of the reported thesis run, not a superseded routing experiment.
% Horizon check: Treat encoder history, scored training window, joint-estimation horizon, and free-run validation horizon separately; relate candidates to relevant stable modes and test sensitivity to slower dynamics and free-integrator accumulation.
% Implementation audit: Read scripts/gantry/gantry_interconnect_dynamic.py, gantry_dynamic/{model,config,controller,data,training}.py, and model_augmentation/fit_systems/{interconnect,closed_loop,pre_encoder}.py.
% Reference check: Verify sources for parallel LFR augmentation, encoder initialization, closed-loop state extension, and any claimed architectural precedent against the original paper or thesis.
% Cautions: Earlier routing plans are superseded; distinguish the truth-only hidden absorber from learned states; closed-loop fit alone does not prove autonomous plant recovery.
% Figures: Absorber schematic, author's updated architecture, and thesis-owned common-reference loops. No residual panel or horizon figure.
% Section is finished when: Every signal has a defined frame and timing, the RK4 boundary is explicit, and the described architecture matches the final code.
%
% SOURCE MAP
% Hoekstra et al. 2025 EJC: dynamic-parallel topology, additional states, encoder-based truncated simulation.
% Hoekstra et al. 2026 Automatica submission: extended LFR formulation, baseline-preserving model/encoder initialization, joint identification context.
% Hoekstra et al. 2026 encoder paper: reconstructability-based W^b initialization; it does not identify W^a or test dynamic augmentation.
% Kessels 2025 thesis: direct closed-loop rollout, controller equations, and reconstruction of the machine controller state for each window.
% Drenth 2025 thesis: LPV-specific augmentation precedent; not the unrestricted MLP structure used here.
% This project: gantry specialization, residual-loop implementation, exact routing, RK4 boundary, and verification.
\section{Dynamic LPV-LFR Augmentation}\label{sec:augmentation}

This section extends dynamic-parallel augmentation to the self-scheduled
baseline of Section~\ref{sec:lfr} and to records measured under feedback.
The physical parameters $\theta_{\mathrm{base}}$, the network parameters
$\theta_{\mathrm{aug}}$ and the encoder parameters $\eta$ are estimated
jointly.
Section~\ref{sec:aug_structure} defines the augmented model,
Section~\ref{sec:aug_encoder} the initial state of each training window, and
Section~\ref{sec:aug_closed_loop} the closed-loop objective.

\subsection{Augmented model}\label{sec:aug_structure}
The baseline omits effects of the machine.
Garc\'ia-Herreros et al.\ neglect the resonance of the supporting base, the
vibration of the cross-arm on its flexible plates, the detent forces of the
linear motors, and the variation of friction along the
stroke~\cite[Sec.~2.4]{garcia2013model}.
The baseline \eqref{eq:eom} also leaves out their Coulomb
friction~\cite[Sec.~2.2]{garcia2013model}.
If the omitted effects are dynamic, the six baseline states cannot represent
them.
Refitting $\theta_{\mathrm{base}}$ adds no states, and neither does a static
correction of the state transition~\cite[Sec.~2]{hoekstra2025lfr}.
The augmentation therefore needs states of its own.
We identify the dynamic-parallel model~\cite[Table~1]{hoekstra2025lfr}
\begin{equation}
\begin{aligned}
 \tilde x_{k+1}
 &=f_{\mathrm{base}}(\theta_{\mathrm{base}},\tilde x_k,u_k)
   +f_{\mathrm{aug}}(\theta_{\mathrm{aug}},\hat x_k,u_k),\\
 \bar x_{k+1}
 &=g_{\mathrm{aug}}(\theta_{\mathrm{aug}},\hat x_k,u_k),\\
 \hat y_k&=h_{\mathrm{base}}(\tilde x_k)=C_n\tilde x_k ,
\end{aligned}
\label{eq:aug_transition}
\end{equation}
where $\tilde x=\col(q,\dot q)\in\R^6$ is the physical state of
Section~\ref{sec:lfr}, $\bar x\in\R^{n_{x_a}}$ collects the additional
states, $\hat x=\col(\tilde x,\bar x)$, $u=P\uact$ is the generalised force
of \eqref{eq:Ptransform}, and $\hat y$ estimates the measured positions
$y=\qact$ (Fig.~\ref{fig:aug_structure}).
We use the parallel structure because the orthogonal construction of
Section~\ref{sec:obc} is formulated for additive
augmentation~\cite{gyorok2026obc}.
Drenth augments a discrete-time LPV-LFR baseline with learning blocks that
are themselves LPV-LFR~\cite[Sec.~5.2]{drenth2025thesis}.
Here, an unrestricted network acts in parallel to a continuous-time
baseline.

We take $f_{\mathrm{base}}$ as one RK4 step of \eqref{eq:eom} over the sample
period $T_s$, with $u$ held over the step.
Each stage evaluates $M(Y)$ at its own state, so $f_{\mathrm{base}}$
remains self-scheduled.
$f_{\mathrm{aug}}$ is added after the step, so it is a discrete-time
correction to physics integrated in continuous time.

All signals are normalised with the statistics of the training
records~\cite[Sec.~5.3]{hoekstra2026lfr}.
Hoekstra et al.\ take the state statistics from a baseline simulation.
We take them from $q=P^{-\top}y$ and its first difference.

One multilayer perceptron with tanh hidden
layers~\cite[Eq.~(14)]{hoekstra2026lfr} supplies both learned functions,
\begin{equation}
 \begin{bmatrix}
  f_{\mathrm{aug}}(\theta_{\mathrm{aug}},\hat x_k,u_k)\\
  g_{\mathrm{aug}}(\theta_{\mathrm{aug}},\hat x_k,u_k)
 \end{bmatrix}
 =\phi_{\mathrm{aug}}\bigl(\theta_{\mathrm{aug}},\col(\hat x_k,u_k)\bigr),
 \label{eq:aug_mlp}
\end{equation}
where $\theta_{\mathrm{aug}}$ collects its weights and biases.
The first six outputs form $f_{\mathrm{aug}}$.
They correct every physical row, the positions included, because the
dynamic-parallel $f_{\mathrm{aug}}$ acts on the whole baseline
state~\cite[Sec.~2]{hoekstra2025lfr}.
A correction restricted to the $\Theta$ rows could not reach the coupled $X$
and $Y$ response directly.
The last $n_{x_a}$ outputs form $g_{\mathrm{aug}}$ and are the next
additional state itself.

\begin{figure}[t]
 \centering
 \resizebox{\columnwidth}{!}{\input{figures/tex/augmentation_structure}}
 \caption{Augmented model \eqref{eq:aug_transition}. The physics step
 $f_{\mathrm{base}}$ and the network $\phi_{\mathrm{aug}}$ act in parallel on
 $\hat x_k$. $\phi_{\mathrm{aug}}$ supplies $f_{\mathrm{aug}}$ and
 $g_{\mathrm{aug}}$, and the output map is not learned
 ($h_{\mathrm{aug}}=0$). The encoder $\psi$ sets $\hat x$ at each window
 start. During training, $u_k$ is the closed-loop input $\hat u_k$ of
 Section~\ref{sec:aug_closed_loop}.}
 \label{fig:aug_structure}
\end{figure}

The output map $C_n$ is the kinematic map
$\begin{bmatrix}P^\top & 0\end{bmatrix}$ of \eqref{eq:Ptransform} in
normalised coordinates, so the positions are read from the physical state
alone.
The additional states reach the output only through $f_{\mathrm{aug}}$, one
step later.
They have no assigned physical meaning and are not the states of an omitted
effect of the simulated machine (Section~\ref{sec:setup}).
The augmentation is therefore judged at the output, not by its states.

No learned function feeds the input of another within a step, which is the
parallel case $D_{zw}=0$ of~\cite[Eq.~(5)]{hoekstra2026lfr}.
The interconnection graph is therefore acyclic.
Since $\phi_{\mathrm{aug}}$ is smooth and $f_{\mathrm{base}}$ is smooth
wherever $M(Y)$ is nonsingular (Section~\ref{sec:params}), the model is well
posed~\cite[Thm.~9, Remark~10]{hoekstra2026lfr}.

Training starts from the baseline, as in~\cite[Eq.~(29)]{hoekstra2026lfr}.
The output layer of $\phi_{\mathrm{aug}}$ is initialised at zero,
\begin{equation}
 W_L=0,\quad b_L=0
 \quad\Longrightarrow\quad
 f_{\mathrm{aug}}=0,\quad g_{\mathrm{aug}}=0 ,
 \label{eq:aug_zero_init}
\end{equation}
where $W_L$ and $b_L$ are the weight and bias of that layer.
The augmented model then reproduces the baseline from the same physical
initial state.
Hoekstra et al.\ use a network with a linear bypass and allow the
additional-state rows to start
nonzero~\cite[Eq.~(29), Sec.~5.4.3]{hoekstra2026lfr}.
We use no bypass, as in their published dynamic-parallel example code, so
$\bar x_{k+1}=0$ at the start.
In closed loop, training therefore starts from the baseline under its
controller, not from a random model, which is prone to closed-loop
instability~\cite[Sec.~5.3.2.2]{kessels2025ai}.
From this start, $f_{\mathrm{aug}}$ can also learn changes that a different
$\theta_{\mathrm{base}}$ would produce, so the split between the two parts
is not unique~\cite[Sec.~5.2]{hoekstra2026lfr}.
Section~\ref{sec:obc} constrains this split.

\subsection{Encoder and initial state}\label{sec:aug_encoder}
Each training window simulates \eqref{eq:aug_transition} from an initial
state, but only the positions are measured.
As in the subspace encoder method~\cite[Sec.~3]{beintema2023subnet}, an
encoder estimates this state from the input and output history,
\begin{equation}
 \hat x_{\tau|\tau}
 =\psi_\eta(\xi_\tau),\qquad
 \xi_\tau=\col\bigl(y_{\tau-n:\tau},\,u_{\tau-n:\tau}\bigr),
 \label{eq:aug_encoder}
\end{equation}
where $\tau$ is the window start and $n$ the encoder history.
Both histories include sample $\tau$, as the reconstructability
map~\cite[Eq.~(15)]{hoekstra2026encoder} does.
The encoder estimates the full state $\hat x$, the measured positions
included.

The encoder is a linear map with a nonlinear
correction~\cite[Eq.~(8)]{hoekstra2026encoder},
\begin{equation}
 \psi_\eta(\xi_\tau)
 =\begin{bmatrix}W^b\\ W^a\end{bmatrix}\xi_\tau+\tilde\psi(\xi_\tau),
 \label{eq:aug_enc_lin}
\end{equation}
where $W^b$ and $W^a$ give the physical and the additional states and
$\tilde\psi$ is a multilayer perceptron.

The physical rows $W^b$ are initialised from the baseline.
For a nonlinear baseline, Hoekstra et al.\ build $W^b$ from a linearisation
about an equilibrium~\cite[Eqs.~(30), (31)]{hoekstra2026encoder}.
We linearise \eqref{eq:eom} about rest at $Y=0$, with the nominal parameters
of Appendix~\ref{app:params}, which gives
\begin{equation}
 A_c=\begin{bmatrix}0 & I\\ -M_0^{-1}K & -M_0^{-1}C\end{bmatrix},\quad
 B_c=\begin{bmatrix}0\\ M_0^{-1}\end{bmatrix},\quad
 C_d=\begin{bmatrix}P^\top & 0\end{bmatrix},
 \label{eq:aug_lin}
\end{equation}
where $M_0=M(0)$ from \eqref{eq:Msplit}.
$(A_d,B_d)$ is the zero-order-hold discretisation of $(A_c,B_c)$ over $T_s$.
$W^b$ is then the reconstructability
map~\cite[Eqs.~(16), (17)]{hoekstra2026encoder}
\begin{equation}
 W^b=\begin{bmatrix}
  A_d^nO_n^{+} & \;r_n-A_d^nO_n^{+}T_n
 \end{bmatrix},
 \label{eq:aug_wb}
\end{equation}
where $O_n$ is the observability matrix of $(A_d,C_d)$ over the history $n$,
$T_n$ the Toeplitz matrix of input-to-output responses, $r_n$ the
input-to-state map, and the two column blocks act on the output and input
parts of $\xi_\tau$.
All matrices are normalised as the signals.
The measured positions make $(A_d,C_d)$ observable, so $O_n$ has full column
rank and $O_n^{+}$ is a left inverse.
The linearisation only sets the starting value of $W^b$, since the encoder
is trained jointly with the model~\cite[Sec.~3.1]{hoekstra2026encoder}.

The rows $W^a$ have no baseline counterpart and are drawn from a Kaiming
uniform distribution, where~\cite[Sec.~5.4.2]{hoekstra2026lfr} uses Xavier.
The correction $\tilde\psi$ starts at zero.
By \eqref{eq:aug_zero_init}, the encoded additional states do not affect the
output at the start of training.

\subsection{Closed-loop rollout and objective}\label{sec:aug_closed_loop}
The records are measured under feedback.
The model is used to predict the machine under feedback, with a known and
possibly changed controller.
We therefore fit the closed-loop response, instead of an open-loop
simulation driven by the recorded input as
in~\cite[Eq.~(22)]{hoekstra2026lfr}.
Kessels trains an augmented model with the known controller inside each
truncated rollout and propagates gradients through
it~\cite[Eqs.~(5.12), (5.13c), (5.13d)]{kessels2025ai}.
We adopt this and keep the controller outside the learned model, so that it
can be replaced in evaluation~\cite[Sec.~5.3.2.3]{kessels2025ai}.

The baseline gives a second reason.
$X$ and $Y$ have no stiffness in \eqref{eq:damping_stiffness_matrices}, so the
baseline integrates both freely for every $Y$.
In an open-loop rollout, a velocity or force error accumulates in these
positions and does not decay.
Inside the loop, the controller acts on such errors as it does on the
machine.

The loop also suppresses model errors within its
bandwidth~\cite[Sec.~5.3.2.3]{kessels2025ai}.
A small closed-loop error therefore does not show that the plant model is
accurate in open loop, which Section~\ref{sec:results} tests separately.

The recorded loop is driven by a reference $r$ and a feedforward
$u^{\mathrm{ff}}$ through a linear controller with state $s$ and matrices
$(A_K,B_K,C_K,D_K)$, which outputs actuator forces.
The model loop receives the same $r$, $u^{\mathrm{ff}}$ and controller, and
differs only in its output $\hat y$ (Fig.~\ref{fig:aug_closed_loop}),
\begin{equation}
\begin{aligned}
 s_{k+1}&=A_K s_k+B_K(r_k-y_k), &
 u_{\mathrm{act},k}&=C_K s_k+D_K(r_k-y_k)+u^{\mathrm{ff}}_k,\\
 \hat s_{k+1}&=A_K \hat s_k+B_K(r_k-\hat y_k), &
 \hat u_{\mathrm{act},k}&=C_K \hat s_k+D_K(r_k-\hat y_k)+u^{\mathrm{ff}}_k .
\end{aligned}
\label{eq:aug_direct_controller}
\end{equation}
Subtracting the recorded loop from the model loop removes $r$ and
$u^{\mathrm{ff}}$, which gives the residual form we implement,
\begin{equation}
\begin{aligned}
 x^c_{k+1}&=A_K x^c_k+B_K(y_k-\hat y_k),\qquad x^c_\tau=0,\\
 \hat u_k&=u_k+P\bigl(C_K x^c_k+D_K(y_k-\hat y_k)\bigr),
\end{aligned}
\label{eq:aug_controller}
\end{equation}
where $x^c=\hat s-s$, $u_k=P u_{\mathrm{act},k}$ is the recorded
generalised force, and $\hat u_k=P\hat u_{\mathrm{act},k}$ drives
\eqref{eq:aug_transition}.
The model thus receives the recorded input plus the controller's response to
its own output error.
This form is exact when both loops use the same linear controller at the
same rate, not scheduled on the model state, with the same feedforward and
without saturation.

\begin{figure}[t]
 \centering
 \includegraphics[width=\columnwidth]{closed_loop_same_reference.pdf}
 \caption{Recorded plant loop (top) and augmented-model loop (bottom),
 driven by one shared reference $r_k$ and feedforward $u_k^{\mathrm{ff}}$,
 with the same feedback controller. The model output is evaluated
 before the input advances its state to the next sample.}
 \label{fig:aug_closed_loop}
\end{figure}

The condition $x^c_\tau=0$ starts the model's controller in the state of the
recorded controller, $\hat s_\tau=s_\tau$, without reconstructing that state.
Kessels instead reconstructs it from the measured output, the reference and
the controller~\cite[Remark~5.4]{kessels2025ai}.
Model error from before $\tau$ is therefore not carried into the window.
The model has no feedthrough, so each sample is evaluated in a fixed order:
$\hat y_k$ from the state, then $\hat u_k$ from \eqref{eq:aug_controller},
then both states advance.

Over the window starts $\mathcal T$ and the window length $n_f$, the
parameters $\vartheta=(\theta_{\mathrm{base}},\theta_{\mathrm{aug}},\eta)$
minimise the output error
\begin{equation}
 V(\vartheta)=\frac{1}{|\mathcal T|\,n_f\,n_y}
 \sum_{\tau\in\mathcal T}\sum_{\ell=0}^{n_f-1}
 \norm{y_{\tau+\ell}-\hat y_{\tau+\ell|\tau}}_2^2
 \label{eq:aug_loss}
\end{equation}
subject to \eqref{eq:aug_transition}, \eqref{eq:aug_encoder} and
\eqref{eq:aug_controller}, where $n_y=3$ and both outputs are normalised.
This is the truncated objective of~\cite[Eq.~(22)]{hoekstra2026lfr} with
each window simulated in closed loop, as
in~\cite[Eq.~(5.12)]{kessels2025ai}.
As in both, every sample of the window is scored.
$\theta_{\mathrm{base}}$ is estimated in the coordinates of
Section~\ref{sec:params}, which keep $M(Y)$ positive definite at every
iterate.
The controller is fixed.
In the loop, a model error reaches the output through a decaying closed-loop
response.
The window length $n_f$ is chosen to span this response.
Validation simulates each validation record as one closed-loop window from
the first sample with a full encoder history, as
in~\cite[Eq.~(5.14)]{kessels2025ai}.
Section~\ref{sec:setup} gives $n_{x_a}$, the network sizes, $n$, $n_f$ and
the optimiser.
